%!TEX TS-program = xelatex
%!TEX encoding = UTF-8 Unicode
\documentclass[10pt]{article}
\title{Java vs Python --- A cheatsheet}
\author{Mathieu N\'eron}
\usepackage{cheatsheet-style}

\begin{document}
\cheatsheetTitle{Java vs Python}{Mathieu N\'eron}{2026}

\noindent\small This cheatsheet anchors Python on constructs you already know in Java.
Left column: Java. Right column: Python (3.10+). Code is illustrative, not always
the most idiomatic possible.

%===============================================================================
\cheatsection{Hello world \& program entry}
%===============================================================================

\begin{construct}{Program entry point}
\noindent
\begin{minipage}[t]{0.485\linewidth}
\begin{lstlisting}[style=java, title=Java]
public class Hello {
  public static void main(String[] args) {
    System.out.println("Hello, " + args[0]);
  }
}
// javac Hello.java && java Hello world
\end{lstlisting}
\end{minipage}\hfill
\begin{minipage}[t]{0.485\linewidth}
\begin{lstlisting}[style=python, title=Python]
import sys

def main():
    print(f"Hello, {sys.argv[1]}")

if __name__ == "__main__":
    main()
# python hello.py world
\end{lstlisting}
\end{minipage}
\end{construct}

\begin{construct}{File / module layout}
\noindent
\begin{minipage}[t]{0.485\linewidth}
\begin{lstlisting}[style=java, title=Java]
// One public class per file.
// Filename must match class name.
// src/main/java/com/acme/Hello.java
package com.acme;
public class Hello { ... }
\end{lstlisting}
\end{minipage}\hfill
\begin{minipage}[t]{0.485\linewidth}
\begin{lstlisting}[style=python, title=Python]
# Any number of top-level defs/classes per file.
# A file IS a module. A directory with
# __init__.py (or PEP 420 implicit) IS a package.
# acme/hello.py     --> import acme.hello
\end{lstlisting}
\end{minipage}
\end{construct}

%===============================================================================
\cheatsection{Variables, constants, type inference}
%===============================================================================

\begin{construct}{Local variables}
\noindent
\begin{minipage}[t]{0.485\linewidth}
\begin{lstlisting}[style=java, title=Java]
int x = 5;                     // primitive
Integer y = 5;                 // boxed
String s = "hi";
var list = new ArrayList<Integer>(); // Java 10+
final double PI = 3.14159;     // immutable binding
\end{lstlisting}
\end{minipage}\hfill
\begin{minipage}[t]{0.485\linewidth}
\begin{lstlisting}[style=python, title=Python]
x = 5                          # dynamic
x: int = 5                     # PEP 526 type hint
s: str = "hi"
nums: list[int] = []           # PEP 585 (3.9+)
PI: Final[float] = 3.14159     # PEP 591
# Convention: ALL_CAPS = "constant" (not enforced)
\end{lstlisting}
\end{minipage}
\end{construct}

\begin{construct}{Scope rules}
\noindent
\begin{minipage}[t]{0.485\linewidth}
\begin{lstlisting}[style=java, title=Java]
// Block-scoped. Shadowing forbidden in same block.
{ int x = 1; }
// 'this.x' for instance, 'ClassName.x' for static.
\end{lstlisting}
\end{minipage}\hfill
\begin{minipage}[t]{0.485\linewidth}
\begin{lstlisting}[style=python, title=Python]
# Function-scoped (LEGB: local, enclosing, global,
# built-in). No block scope inside if/for/while.
def f():
    global g_count       # rebind module-level
    nonlocal n           # rebind enclosing fn
\end{lstlisting}
\end{minipage}
\end{construct}

%===============================================================================
\cheatsection{Primitive types \& conversions}
%===============================================================================

\begin{construct}{Numeric types}
\noindent
\begin{minipage}[t]{0.485\linewidth}
\begin{lstlisting}[style=java, title=Java]
byte b; short s; int i; long l;     // 8/16/32/64-bit
float f; double d;                  // 32/64-bit IEEE
int n = Integer.parseInt("42");
double x = Double.parseDouble("1.5");
String t = String.valueOf(42);
// Overflow wraps silently (use Math.addExact for exc.)
\end{lstlisting}
\end{minipage}\hfill
\begin{minipage}[t]{0.485\linewidth}
\begin{lstlisting}[style=python, title=Python]
# int  -> arbitrary precision (no overflow)
# float -> 64-bit IEEE
# complex -> a+bj, bool subclasses int
n = int("42")
x = float("1.5")
t = str(42)
big = 2 ** 200          # works fine
\end{lstlisting}
\end{minipage}
\end{construct}

\begin{construct}{Bool, char, byte}
\noindent
\begin{minipage}[t]{0.485\linewidth}
\begin{lstlisting}[style=java, title=Java]
boolean b = true;
char c = 'A';                       // 16-bit UTF-16 unit
byte[] bytes = "hi".getBytes(StandardCharsets.UTF_8);
\end{lstlisting}
\end{minipage}\hfill
\begin{minipage}[t]{0.485\linewidth}
\begin{lstlisting}[style=python, title=Python]
b = True                            # True / False
c = "A"                             # no char type
bs: bytes = "hi".encode("utf-8")    # bytes literal: b"hi"
ba: bytearray = bytearray(b"hi")    # mutable
\end{lstlisting}
\end{minipage}
\end{construct}

%===============================================================================
\cheatsection{Strings}
%===============================================================================

\begin{construct}{Interpolation \& formatting}
\noindent
\begin{minipage}[t]{0.485\linewidth}
\begin{lstlisting}[style=java, title=Java]
String name = "world";
String s1 = "Hello, " + name;
String s2 = String.format("%s = %.2f", "pi", 3.14159);
// Java 15+: text blocks (multi-line)
String body = """
    Hello,
    World
    """;
\end{lstlisting}
\end{minipage}\hfill
\begin{minipage}[t]{0.485\linewidth}
\begin{lstlisting}[style=python, title=Python]
name = "world"
s1 = f"Hello, {name}"               # PEP 498 f-string
s2 = "{} = {:.2f}".format("pi", 3.14159)
s3 = "%s = %.2f" % ("pi", 3.14159)  # legacy
body = """
    Hello,
    World
    """                             # multi-line literal
\end{lstlisting}
\end{minipage}
\end{construct}

\begin{construct}{Common operations}
\noindent
\begin{minipage}[t]{0.485\linewidth}
\begin{lstlisting}[style=java, title=Java]
s.length(); s.toUpperCase(); s.trim();
s.startsWith("H"); s.contains("ll");
String[] parts = s.split(",");
String joined = String.join("-", parts);
// Strings IMMUTABLE; use StringBuilder for builds
\end{lstlisting}
\end{minipage}\hfill
\begin{minipage}[t]{0.485\linewidth}
\begin{lstlisting}[style=python, title=Python]
len(s); s.upper(); s.strip()
s.startswith("H"); "ll" in s
parts = s.split(",")
joined = "-".join(parts)
# Strings IMMUTABLE; "".join(list) is the build idiom
\end{lstlisting}
\end{minipage}
\end{construct}

%===============================================================================
\cheatsection{Control flow}
%===============================================================================

\begin{construct}{If / else / ternary}
\noindent
\begin{minipage}[t]{0.485\linewidth}
\begin{lstlisting}[style=java, title=Java]
if (x > 0) { ... }
else if (x == 0) { ... }
else { ... }

int sign = x > 0 ? 1 : (x < 0 ? -1 : 0);
\end{lstlisting}
\end{minipage}\hfill
\begin{minipage}[t]{0.485\linewidth}
\begin{lstlisting}[style=python, title=Python]
if x > 0:
    ...
elif x == 0:
    ...
else:
    ...

sign = 1 if x > 0 else (-1 if x < 0 else 0)
\end{lstlisting}
\end{minipage}
\end{construct}

\begin{construct}{Switch / match}
\noindent
\begin{minipage}[t]{0.485\linewidth}
\begin{lstlisting}[style=java, title=Java]
// Java 14+ switch expression with pattern matching (Java 21)
String label = switch (shape) {
  case Circle c       -> "circle r=" + c.r();
  case Rectangle(var w, var h) -> "rect " + w + "x" + h;
  case null           -> "none";
  default             -> "unknown";
};
\end{lstlisting}
\end{minipage}\hfill
\begin{minipage}[t]{0.485\linewidth}
\begin{lstlisting}[style=python, title=Python]
# 3.10+ structural pattern matching (PEP 634)
match shape:
    case Circle(r=r):
        label = f"circle r={r}"
    case Rectangle(w=w, h=h):
        label = f"rect {w}x{h}"
    case None:
        label = "none"
    case _:
        label = "unknown"
\end{lstlisting}
\end{minipage}
\end{construct}

%===============================================================================
\cheatsection{Loops}
%===============================================================================

\begin{construct}{Counting / classic for}
\noindent
\begin{minipage}[t]{0.485\linewidth}
\begin{lstlisting}[style=java, title=Java]
for (int i = 0; i < 10; i++) { ... }
for (int i = 9; i >= 0; i--) { ... }
\end{lstlisting}
\end{minipage}\hfill
\begin{minipage}[t]{0.485\linewidth}
\begin{lstlisting}[style=python, title=Python]
for i in range(10):           # 0..9
    ...
for i in range(9, -1, -1):    # 9..0
    ...
\end{lstlisting}
\end{minipage}
\end{construct}

\begin{construct}{For-each / collection iteration}
\noindent
\begin{minipage}[t]{0.485\linewidth}
\begin{lstlisting}[style=java, title=Java]
for (String name : names) { ... }
for (var e : map.entrySet()) {
  System.out.println(e.getKey() + "=" + e.getValue());
}
\end{lstlisting}
\end{minipage}\hfill
\begin{minipage}[t]{0.485\linewidth}
\begin{lstlisting}[style=python, title=Python]
for name in names:
    ...
for k, v in d.items():        # tuple unpacking
    print(f"{k}={v}")
for i, x in enumerate(xs):    # index + value
    ...
\end{lstlisting}
\end{minipage}
\end{construct}

\begin{construct}{While / do-while}
\noindent
\begin{minipage}[t]{0.485\linewidth}
\begin{lstlisting}[style=java, title=Java]
while (cond) { ... }
do { ... } while (cond);
break; continue;
outer: for (...) { for (...) { break outer; } }
\end{lstlisting}
\end{minipage}\hfill
\begin{minipage}[t]{0.485\linewidth}
\begin{lstlisting}[style=python, title=Python]
while cond:
    ...
# No do-while; idiom:
while True:
    ...
    if not cond: break
break; continue
# No labeled break; refactor or use a flag/return
\end{lstlisting}
\end{minipage}
\end{construct}

%===============================================================================
\cheatsection{Collections \& data structures}
%===============================================================================

\cheatsub{Quick reference}

\noindent
\begin{tabularx}{\linewidth}{@{}>{\ttfamily\small}l >{\small}l X@{}}
\toprule
\rowcolor{accentDark}
\refhead{\normalfont Java} & \refhead{Behavior} & \refhead{Python} \\
\midrule
ArrayList<T>       & dynamic array, O(1) get & \code{list} \\
LinkedList<T>      & doubly-linked, O(1) ends & \code{collections.deque} \\
int[] / Object[]   & fixed-size primitive array & \code{array.array("i", ...)} (typed) \\
HashSet<T>         & unordered set, O(1) ops & \code{set} \\
LinkedHashSet<T>   & insertion-ordered set & \code{dict.fromkeys(xs)} (3.7+ keeps order) \\
TreeSet<T>         & sorted set & \code{sortedcontainers.SortedSet} (PyPI) \\
EnumSet            & bit-set over enum & \code{set} of \code{Enum} members \\
HashMap<K,V>       & unordered map & \code{dict} (insertion-ordered since 3.7!) \\
LinkedHashMap<K,V> & insertion-ordered map & \code{dict} (or \code{collections.OrderedDict} for \code{move\_to\_end}) \\
TreeMap<K,V>       & sorted map & \code{sortedcontainers.SortedDict} (PyPI) \\
IdentityHashMap    & identity-keyed map & \code{\{id(k): v\}} or \code{weakref.WeakKeyDictionary} \\
EnumMap            & map keyed by enum & plain \code{dict} keyed on \code{Enum} \\
WeakHashMap        & weakly-referenced keys & \code{weakref.WeakKeyDictionary} \\
ArrayDeque (queue) & FIFO queue, ring buffer & \code{collections.deque} \\
ArrayDeque (stack) & LIFO stack & \code{list} via \code{append/pop} \\
ArrayDeque         & double-ended queue & \code{collections.deque} \\
PriorityQueue<T>   & binary heap (min-heap) & \code{heapq} on a list \\
Stack<T>           & legacy LIFO (synchronized) & \code{list} via \code{append/pop} \\
ConcurrentHashMap  & thread-safe map & \code{dict} + \code{threading.Lock} (or shared via \code{multiprocessing.Manager}) \\
BlockingQueue      & thread-safe FIFO + blocking & \code{queue.Queue} \\
LinkedBlockingDeque & blocking deque & \code{queue.Queue} (no native deque variant) \\
DelayQueue         & expiry-ordered queue & \code{sched} module / hand-roll on \code{heapq} \\
List.copyOf, Map.of & shallow-immutable views & \code{tuple}, \code{frozenset}, \code{MappingProxyType} \\
\bottomrule
\end{tabularx}

\bigskip

\cheatsub{Lists}

\begin{construct}{ArrayList / LinkedList / array}
\noindent
\begin{minipage}[t]{0.485\linewidth}
\begin{lstlisting}[style=java, title=Java]
List<Integer> a = new ArrayList<>();         // O(1) get
a.add(1); a.get(0); a.remove(0); a.size();

List<Integer> ll = new LinkedList<>();       // O(1) ends
ll.addFirst(0); ll.addLast(2);

int[] arr = new int[10];                     // fixed
List<Integer> immut = List.of(1, 2, 3);
\end{lstlisting}
\end{minipage}\hfill
\begin{minipage}[t]{0.485\linewidth}
\begin{lstlisting}[style=python, title=Python]
xs = [1, 2, 3]                # dynamic array
xs.append(4); xs[0]; xs.pop(0); len(xs)
xs[1:3]; xs[::-1]             # slice / reverse view-copy

from collections import deque
dq = deque([1,2,3])           # O(1) at both ends
dq.appendleft(0); dq.append(4)

import array
a = array.array("i", [1,2,3]) # typed C-contiguous
t = (1, 2, 3)                 # immutable tuple
\end{lstlisting}
\end{minipage}
\end{construct}

\cheatsub{Sets}

\begin{construct}{HashSet / LinkedHashSet / TreeSet}
\noindent
\begin{minipage}[t]{0.485\linewidth}
\begin{lstlisting}[style=java, title=Java]
Set<String> hs = new HashSet<>();             // unordered
Set<String> ls = new LinkedHashSet<>();       // ins-order
Set<String> ts = new TreeSet<>();             // sorted
ts.first(); ts.headSet("m"); ts.subSet("a","z");

EnumSet<Day> weekend = EnumSet.of(SAT, SUN);
\end{lstlisting}
\end{minipage}\hfill
\begin{minipage}[t]{0.485\linewidth}
\begin{lstlisting}[style=python, title=Python]
s = {"a", "b"}; s.add("c"); s.discard("a")
s | t; s & t; s - t; s ^ t        # ops
fs = frozenset(s)                  # immutable

# Insertion order: dict keys preserve, set does NOT.
ordered = list(dict.fromkeys(["b","a","b","c"]))

# Sorted set: 3rd-party
from sortedcontainers import SortedSet
ts = SortedSet(["c","a","b"]);  ts[0]; ts.irange("a","m")
\end{lstlisting}
\end{minipage}
\end{construct}

\cheatsub{Maps}

\begin{construct}{HashMap / LinkedHashMap / TreeMap}
\noindent
\begin{minipage}[t]{0.485\linewidth}
\begin{lstlisting}[style=java, title=Java]
Map<String,Integer> hm = new HashMap<>();        // unord
Map<String,Integer> lm = new LinkedHashMap<>();  // ins-ord
Map<String,Integer> tm = new TreeMap<>();        // sorted

hm.put("a", 1); hm.getOrDefault("b", 0);
hm.computeIfAbsent("c", k -> 42);
hm.merge("a", 1, Integer::sum);                  // tally

tm.firstKey(); tm.headMap("m"); tm.floorKey("g");
\end{lstlisting}
\end{minipage}\hfill
\begin{minipage}[t]{0.485\linewidth}
\begin{lstlisting}[style=python, title=Python]
# dict preserves INSERTION ORDER (3.7+). So plain dict
# already covers HashMap AND LinkedHashMap.
d = {"a": 1}
d.get("b", 0)                              # getOrDefault
d.setdefault("c", 42)                      # computeIfAbsent

# computeIfAbsent for collection-valued maps:
from collections import defaultdict
groups = defaultdict(list); groups["k"].append("v")

# Tally / merge:
from collections import Counter
c = Counter("aabbbccd"); c.most_common(2)

# Sorted map: 3rd-party
from sortedcontainers import SortedDict
sm = SortedDict({"b":2,"a":1}); sm.irange("a","m")
\end{lstlisting}
\end{minipage}
\end{construct}

\begin{construct}{Identity / weak / read-only / enum keys}
\noindent
\begin{minipage}[t]{0.485\linewidth}
\begin{lstlisting}[style=java, title=Java]
Map<Key,V> id = new IdentityHashMap<>();        // == not equals
Map<Key,V> wk = new WeakHashMap<>();            // GC-able keys
Map<K,V>   ro = Collections.unmodifiableMap(m); // view
Map<Day,V> em = new EnumMap<>(Day.class);
\end{lstlisting}
\end{minipage}\hfill
\begin{minipage}[t]{0.485\linewidth}
\begin{lstlisting}[style=python, title=Python]
import weakref
wk = weakref.WeakKeyDictionary()           # WeakHashMap
wv = weakref.WeakValueDictionary()         # values can vanish

from types import MappingProxyType
ro = MappingProxyType({"a": 1})            # read-only view

# IdentityHashMap idiom:
id_map = {}
id_map[id(obj)] = value

# EnumMap analog: just dict[Enum, V]
from enum import Enum
class Day(Enum): MON=1; TUE=2
schedule: dict[Day, str] = {}
\end{lstlisting}
\end{minipage}
\end{construct}

\cheatsub{Queue, deque, stack, priority queue}

\begin{construct}{Queue (FIFO) \& Deque}
\noindent
\begin{minipage}[t]{0.485\linewidth}
\begin{lstlisting}[style=java, title=Java]
Deque<Integer> q = new ArrayDeque<>();   // recommended FIFO
q.offer(1); q.offer(2);
int head = q.poll();                     // null if empty
int peek = q.peek();

// Double-ended:
q.offerFirst(0); q.offerLast(3);
q.pollFirst(); q.pollLast();

// Capacity-bounded blocking:
BlockingQueue<Integer> bq = new ArrayBlockingQueue<>(8);
bq.put(42); int v = bq.take();           // blocks
\end{lstlisting}
\end{minipage}\hfill
\begin{minipage}[t]{0.485\linewidth}
\begin{lstlisting}[style=python, title=Python]
from collections import deque
dq = deque(maxlen=8)            # optional cap (drops old)
dq.append(1); dq.appendleft(0)  # O(1) both ends
dq.popleft()                    # FIFO take
dq.pop()                        # LIFO take
dq[0]                           # peek

# Thread-safe blocking queue:
import queue
bq = queue.Queue(maxsize=8)
bq.put(42)              # blocks if full
v = bq.get()            # blocks if empty
# Variants: queue.LifoQueue, queue.PriorityQueue
\end{lstlisting}
\end{minipage}
\end{construct}

\begin{construct}{Stack (LIFO)}
\noindent
\begin{minipage}[t]{0.485\linewidth}
\begin{lstlisting}[style=java, title=Java]
// Avoid legacy java.util.Stack (synchronized + Vector).
Deque<Integer> st = new ArrayDeque<>();
st.push(1); st.push(2);
int top = st.peek();
int v   = st.pop();
\end{lstlisting}
\end{minipage}\hfill
\begin{minipage}[t]{0.485\linewidth}
\begin{lstlisting}[style=python, title=Python]
# Idiomatic: just a list with append/pop.
st = []
st.append(1); st.append(2)
top = st[-1]
v   = st.pop()        # O(1) at the end
# (Use deque if you need both ends; list is faster for LIFO.)
\end{lstlisting}
\end{minipage}
\end{construct}

\begin{construct}{PriorityQueue (heap)}
\noindent
\begin{minipage}[t]{0.485\linewidth}
\begin{lstlisting}[style=java, title=Java]
PriorityQueue<Integer> pq = new PriorityQueue<>();   // min
PriorityQueue<Job> jobs =
    new PriorityQueue<>(Comparator.comparingInt(Job::priority));
pq.offer(3); pq.offer(1); pq.offer(2);
int min = pq.poll();           // 1
\end{lstlisting}
\end{minipage}\hfill
\begin{minipage}[t]{0.485\linewidth}
\begin{lstlisting}[style=python, title=Python]
import heapq
h = []
heapq.heappush(h, 3); heapq.heappush(h, 1); heapq.heappush(h, 2)
mn = heapq.heappop(h)          # 1 (min-heap)

# Custom priority: push (priority, item) tuples
heapq.heappush(h, (job.priority, job))

# Max-heap: push negated keys, or use a wrapper.
# Thread-safe: queue.PriorityQueue.
\end{lstlisting}
\end{minipage}
\end{construct}

\cheatsub{Concurrent \& specialized}

\begin{construct}{Concurrent maps and queues}
\noindent
\begin{minipage}[t]{0.485\linewidth}
\begin{lstlisting}[style=java, title=Java]
ConcurrentHashMap<String,Long> chm = new ConcurrentHashMap<>();
chm.merge("k", 1L, Long::sum);             // atomic tally
chm.computeIfAbsent("k", k -> heavy());     // single-flight

BlockingQueue<Task> q = new LinkedBlockingQueue<>();
q.put(t); Task next = q.take();

CopyOnWriteArrayList<Listener> listeners
    = new CopyOnWriteArrayList<>();         // read-mostly
\end{lstlisting}
\end{minipage}\hfill
\begin{minipage}[t]{0.485\linewidth}
\begin{lstlisting}[style=python, title=Python]
# No native lock-free map. Patterns:
import threading
lock = threading.Lock()
counts = {}
with lock:
    counts["k"] = counts.get("k", 0) + 1

# Thread-safe queue:
import queue
q = queue.Queue()
q.put(t); next_ = q.get()

# Cross-process state:
import multiprocessing
m = multiprocessing.Manager().dict()
\end{lstlisting}
\end{minipage}
\end{construct}

\begin{construct}{Specialized: Counter, ChainMap, NamedTuple}
\noindent
\begin{minipage}[t]{0.485\linewidth}
\begin{lstlisting}[style=java, title=Java]
// Tally: Map<K,Long> + merge, or Stream.collect(groupingBy)
Map<String,Long> tally = words.stream()
    .collect(Collectors.groupingBy(w -> w, Collectors.counting()));

// Layered config: walk a list of Maps yourself.
// Named tuple: record (Java 16+)
record Point(int x, int y) {}
\end{lstlisting}
\end{minipage}\hfill
\begin{minipage}[t]{0.485\linewidth}
\begin{lstlisting}[style=python, title=Python]
from collections import Counter, ChainMap, namedtuple, OrderedDict
c = Counter(words)               # multiset / tally
c["x"] += 1; c.most_common(3)

cfg = ChainMap(env, file_cfg, defaults)   # layered lookup

P = namedtuple("Point", "x y")            # immutable record
# Modern: dataclasses.dataclass(frozen=True) is preferred.

od = OrderedDict()                # adds .move_to_end / popitem(last)
od["a"] = 1; od.move_to_end("a")
\end{lstlisting}
\end{minipage}
\end{construct}

\begin{construct}{Mutability defaults}
\noindent
\begin{minipage}[t]{\linewidth}
\small
\textbf{Java}: \code{ArrayList}, \code{HashMap}, \code{HashSet} mutable.
\code{List.of} / \code{Set.of} / \code{Map.of} return true-immutable instances.
\code{Collections.unmodifiableX} wraps a \emph{view} --- the underlying can still
be mutated by holders of the original.\par
\textbf{Python}: \code{list}, \code{dict}, \code{set} mutable.
\code{tuple}, \code{frozenset}, \code{bytes}, \code{str} immutable.
\code{MappingProxyType} wraps a dict for a shallow-readonly view.
For "deeply readonly" data: \code{@dataclass(frozen=True)},
\code{NamedTuple}, or \code{tuple}-of-tuples.
\end{minipage}
\end{construct}

%===============================================================================
\cheatsection{Functions}
%===============================================================================

\begin{construct}{Definition, default args, varargs}
\noindent
\begin{minipage}[t]{0.485\linewidth}
\begin{lstlisting}[style=java, title=Java]
int add(int a, int b) { return a + b; }
// no default args: use overloading
int add(int a) { return add(a, 1); }
// varargs:
int sum(int... xs) {
  int t = 0; for (int x : xs) t += x; return t;
}
sum(1, 2, 3);
\end{lstlisting}
\end{minipage}\hfill
\begin{minipage}[t]{0.485\linewidth}
\begin{lstlisting}[style=python, title=Python]
def add(a: int, b: int = 1) -> int:
    return a + b

def sum_(*xs: int) -> int:
    return sum(xs)

def cfg(**kw):           # kwargs dict
    print(kw)

add(2)            # b defaults to 1
sum_(1, 2, 3)
cfg(host="x", port=80)
\end{lstlisting}
\end{minipage}
\end{construct}

\begin{construct}{Multiple return values}
\noindent
\begin{minipage}[t]{0.485\linewidth}
\begin{lstlisting}[style=java, title=Java]
// Records (Java 16+) for ad-hoc tuples.
record Pair<A,B>(A first, B second) {}
Pair<Integer,String> divmod(int x, int y) {
  return new Pair<>(x / y, x + " over " + y);
}
\end{lstlisting}
\end{minipage}\hfill
\begin{minipage}[t]{0.485\linewidth}
\begin{lstlisting}[style=python, title=Python]
def divmod_(x: int, y: int) -> tuple[int, int]:
    return x // y, x % y

q, r = divmod_(10, 3)    # tuple unpacking
\end{lstlisting}
\end{minipage}
\end{construct}

%===============================================================================
\cheatsection{Lambdas, closures, higher-order}
%===============================================================================

\begin{construct}{Lambda \& method reference}
\noindent
\begin{minipage}[t]{0.485\linewidth}
\begin{lstlisting}[style=java, title=Java]
Function<Integer,Integer> sq = x -> x * x;
Predicate<String> nonEmpty = s -> !s.isEmpty();
BiFunction<Integer,Integer,Integer> add = (a,b) -> a+b;
// Method ref:
list.forEach(System.out::println);
\end{lstlisting}
\end{minipage}\hfill
\begin{minipage}[t]{0.485\linewidth}
\begin{lstlisting}[style=python, title=Python]
sq = lambda x: x * x         # one-expression only
non_empty = lambda s: bool(s)
add = lambda a, b: a + b

# More common: def + first-class function passing
for x in xs: print(x)        # no need for forEach
\end{lstlisting}
\end{minipage}
\end{construct}

\begin{construct}{Streams / functional pipelines}
\noindent
\begin{minipage}[t]{0.485\linewidth}
\begin{lstlisting}[style=java, title=Java]
List<Integer> result = xs.stream()
    .filter(x -> x > 0)
    .map(x -> x * 2)
    .sorted()
    .toList();
int total = xs.stream().mapToInt(Integer::intValue).sum();
\end{lstlisting}
\end{minipage}\hfill
\begin{minipage}[t]{0.485\linewidth}
\begin{lstlisting}[style=python, title=Python]
result = sorted(x * 2 for x in xs if x > 0)
total = sum(xs)

# Or via map/filter (lazy iterators):
list(map(lambda x: x*2, filter(lambda x: x>0, xs)))
# Idiomatic Python prefers comprehensions.
\end{lstlisting}
\end{minipage}
\end{construct}

%===============================================================================
\cheatsection{Classes}
%===============================================================================

\begin{construct}{Class with fields, constructor, methods}
\noindent
\begin{minipage}[t]{0.485\linewidth}
\begin{lstlisting}[style=java, title=Java]
public class Point {
  private final double x, y;
  public Point(double x, double y) {
    this.x = x; this.y = y;
  }
  public double dist(Point o) {
    double dx = x - o.x, dy = y - o.y;
    return Math.sqrt(dx*dx + dy*dy);
  }
}
\end{lstlisting}
\end{minipage}\hfill
\begin{minipage}[t]{0.485\linewidth}
\begin{lstlisting}[style=python, title=Python]
class Point:
    def __init__(self, x: float, y: float):
        self.x = x
        self.y = y

    def dist(self, o: "Point") -> float:
        dx, dy = self.x - o.x, self.y - o.y
        return (dx * dx + dy * dy) ** 0.5
\end{lstlisting}
\end{minipage}
\end{construct}

\begin{construct}{Records / data classes}
\noindent
\begin{minipage}[t]{0.485\linewidth}
\begin{lstlisting}[style=java, title=Java]
// Java 16+
public record Point(double x, double y) {
  public double dist(Point o) {
    return Math.hypot(x - o.x, y - o.y);
  }
}
// Auto: equals, hashCode, toString, accessors
\end{lstlisting}
\end{minipage}\hfill
\begin{minipage}[t]{0.485\linewidth}
\begin{lstlisting}[style=python, title=Python]
from dataclasses import dataclass

@dataclass(frozen=True, slots=True)
class Point:
    x: float
    y: float
    def dist(self, o: "Point") -> float:
        return ((self.x-o.x)**2 + (self.y-o.y)**2)**0.5
# Auto: __init__, __repr__, __eq__, __hash__
\end{lstlisting}
\end{minipage}
\end{construct}

\begin{construct}{Inheritance, abstract, interfaces}
\noindent
\begin{minipage}[t]{0.485\linewidth}
\begin{lstlisting}[style=java, title=Java]
abstract class Shape {
  abstract double area();
  public String name() { return "Shape"; }
}
interface Drawable { void draw(); }
class Circle extends Shape implements Drawable {
  private final double r;
  public Circle(double r) { this.r = r; }
  @Override double area() { return Math.PI*r*r; }
  @Override public void draw() { ... }
}
\end{lstlisting}
\end{minipage}\hfill
\begin{minipage}[t]{0.485\linewidth}
\begin{lstlisting}[style=python, title=Python]
from abc import ABC, abstractmethod

class Shape(ABC):
    @abstractmethod
    def area(self) -> float: ...
    def name(self) -> str: return "Shape"

class Drawable(ABC):
    @abstractmethod
    def draw(self) -> None: ...

class Circle(Shape, Drawable):     # multiple inherit OK
    def __init__(self, r: float): self.r = r
    def area(self): return 3.14159 * self.r * self.r
    def draw(self): ...
\end{lstlisting}
\end{minipage}
\end{construct}

\begin{construct}{Static / class methods}
\noindent
\begin{minipage}[t]{0.485\linewidth}
\begin{lstlisting}[style=java, title=Java]
public class MathU {
  public static int sq(int x) { return x*x; }
  public static final double PI = 3.14159;
}
MathU.sq(3);
\end{lstlisting}
\end{minipage}\hfill
\begin{minipage}[t]{0.485\linewidth}
\begin{lstlisting}[style=python, title=Python]
class MathU:
    PI = 3.14159

    @staticmethod
    def sq(x: int) -> int: return x * x

    @classmethod
    def named(cls, x): return cls.sq(x)

MathU.sq(3)
\end{lstlisting}
\end{minipage}
\end{construct}

%===============================================================================
\cheatsection{Generics / type parameters}
%===============================================================================

\begin{construct}{Generic class \& method}
\noindent
\begin{minipage}[t]{0.485\linewidth}
\begin{lstlisting}[style=java, title=Java]
class Box<T> {
  private final T value;
  public Box(T v) { this.value = v; }
  public T get() { return value; }
}
<T extends Comparable<T>> T max(T a, T b) {
  return a.compareTo(b) >= 0 ? a : b;
}
List<? extends Number> ns;     // bounded wildcard
\end{lstlisting}
\end{minipage}\hfill
\begin{minipage}[t]{0.485\linewidth}
\begin{lstlisting}[style=python, title=Python]
# 3.12+ syntax (PEP 695)
class Box[T]:
    def __init__(self, v: T): self.value = v
    def get(self) -> T: return self.value

def max_[T: "SupportsRichComparison"](a: T, b: T) -> T:
    return a if a >= b else b

# Pre-3.12: TypeVar from typing
from typing import TypeVar
T = TypeVar("T")
\end{lstlisting}
\end{minipage}
\end{construct}

%===============================================================================
\cheatsection{Error handling}
%===============================================================================

\begin{construct}{Try / catch / finally}
\noindent
\begin{minipage}[t]{0.485\linewidth}
\begin{lstlisting}[style=java, title=Java]
try (var f = Files.newBufferedReader(p)) {
    return f.readLine();
} catch (IOException | RuntimeException e) {
    log.error("read failed", e);
    throw new MyException("oops", e);
} finally {
    cleanup();
}
// Checked exceptions force `throws` declarations.
\end{lstlisting}
\end{minipage}\hfill
\begin{minipage}[t]{0.485\linewidth}
\begin{lstlisting}[style=python, title=Python]
try:
    with open(p) as f:               # context manager
        return f.readline()
except (OSError, RuntimeError) as e:
    log.error("read failed", exc_info=e)
    raise MyError("oops") from e
finally:
    cleanup()
# All exceptions unchecked. No `throws` clause.
\end{lstlisting}
\end{minipage}
\end{construct}

\begin{construct}{Custom exceptions}
\noindent
\begin{minipage}[t]{0.485\linewidth}
\begin{lstlisting}[style=java, title=Java]
public class NotFoundException extends RuntimeException {
  public NotFoundException(String msg) { super(msg); }
}
throw new NotFoundException("user " + id);
\end{lstlisting}
\end{minipage}\hfill
\begin{minipage}[t]{0.485\linewidth}
\begin{lstlisting}[style=python, title=Python]
class NotFoundError(Exception):
    pass

raise NotFoundError(f"user {id}")
\end{lstlisting}
\end{minipage}
\end{construct}

%===============================================================================
\cheatsection{Modules, packages, imports}
%===============================================================================

\begin{construct}{Imports}
\noindent
\begin{minipage}[t]{0.485\linewidth}
\begin{lstlisting}[style=java, title=Java]
package com.acme.svc;
import java.util.List;
import java.util.stream.Collectors;
import static java.util.Map.entry;

// Maven / Gradle: declare in pom.xml / build.gradle
\end{lstlisting}
\end{minipage}\hfill
\begin{minipage}[t]{0.485\linewidth}
\begin{lstlisting}[style=python, title=Python]
from typing import List
from collections import defaultdict
import json
import numpy as np                     # alias

# Pip / Poetry / uv: declare in pyproject.toml
# pip install requests  /  uv add requests
\end{lstlisting}
\end{minipage}
\end{construct}

%===============================================================================
\cheatsection{Concurrency}
%===============================================================================

\begin{construct}{Threads}
\noindent
\begin{minipage}[t]{0.485\linewidth}
\begin{lstlisting}[style=java, title=Java]
Thread t = new Thread(() -> work());
t.start(); t.join();

// Pools:
ExecutorService pool = Executors.newFixedThreadPool(4);
Future<Integer> f = pool.submit(() -> compute());
int v = f.get();
pool.shutdown();
\end{lstlisting}
\end{minipage}\hfill
\begin{minipage}[t]{0.485\linewidth}
\begin{lstlisting}[style=python, title=Python]
import threading
t = threading.Thread(target=work); t.start(); t.join()

# Pools (CPU-bound -> processes; I/O -> threads):
from concurrent.futures import ThreadPoolExecutor
with ThreadPoolExecutor(4) as pool:
    f = pool.submit(compute); v = f.result()
# GIL: pure-Python CPU work doesn't parallelize across
# threads; use multiprocessing or release GIL via C ext.
\end{lstlisting}
\end{minipage}
\end{construct}

\begin{construct}{Async / non-blocking}
\noindent
\begin{minipage}[t]{0.485\linewidth}
\begin{lstlisting}[style=java, title=Java]
CompletableFuture<String> fut =
  CompletableFuture.supplyAsync(() -> fetch())
    .thenApply(String::trim)
    .exceptionally(e -> "fallback");
String r = fut.get();
\end{lstlisting}
\end{minipage}\hfill
\begin{minipage}[t]{0.485\linewidth}
\begin{lstlisting}[style=python, title=Python]
import asyncio

async def run():
    try:
        r = (await fetch()).strip()
    except Exception:
        r = "fallback"
    return r

asyncio.run(run())          # entry point
# async def + await; coroutines run on an event loop.
\end{lstlisting}
\end{minipage}
\end{construct}

%===============================================================================
\cheatsection{Idiom appendix}
%===============================================================================

\begin{construct}{Null / None / equality}
\noindent
\begin{minipage}[t]{0.485\linewidth}
\begin{lstlisting}[style=java, title=Java]
if (s == null) { ... }              // identity
if (a.equals(b)) { ... }            // value
Optional<String> os = Optional.ofNullable(s);
String x = os.orElse("default");
\end{lstlisting}
\end{minipage}\hfill
\begin{minipage}[t]{0.485\linewidth}
\begin{lstlisting}[style=python, title=Python]
if s is None: ...                   # identity (always 'is')
if a == b: ...                      # __eq__
x = s if s is not None else "default"
# Optional[T] in typing == Union[T, None] == T | None (3.10+)
\end{lstlisting}
\end{minipage}
\end{construct}

\begin{construct}{Printing \& logging}
\noindent
\begin{minipage}[t]{0.485\linewidth}
\begin{lstlisting}[style=java, title=Java]
System.out.println("hi");
System.err.printf("x=%d%n", x);
Logger log = LoggerFactory.getLogger(C.class);
log.info("user {} logged in", uid);
\end{lstlisting}
\end{minipage}\hfill
\begin{minipage}[t]{0.485\linewidth}
\begin{lstlisting}[style=python, title=Python]
print("hi")
print(f"x={x}", file=sys.stderr)
import logging
log = logging.getLogger(__name__)
log.info("user %s logged in", uid)
\end{lstlisting}
\end{minipage}
\end{construct}

\begin{construct}{Common gotchas (Python coming from Java)}
\noindent
\begin{minipage}[t]{\linewidth}
\small
\begin{itemize}[leftmargin=*, itemsep=1pt, topsep=1pt]
  \item \textbf{Default arg values are evaluated once.} \code{def f(x=[])} reuses the same list across calls. Use \code{None} sentinels: \code{def f(x=None): x = x or []}.
  \item \textbf{Indentation is syntax.} Mixing tabs and spaces is a real bug. Configure your editor.
  \item \textbf{Slicing copies.} \code{xs[:]} is a shallow copy; \code{xs[a:b] = []} mutates in place.
  \item \textbf{No method overloading.} Latest \code{def} wins. Use default args or \code{functools.singledispatch}.
  \item \textbf{Truthiness is broad.} \code{0}, \code{""}, \code{[]}, \code{\{\}}, \code{None} are all falsy. Prefer \code{is None} for null checks.
  \item \textbf{Equality.} \code{==} calls \code{__eq__} (like Java's \code{equals}). \code{is} is identity (like Java's \code{==}).
  \item \textbf{No private/protected.} Convention: leading \code{\_} = internal, \code{\_\_name} = name-mangled.
  \item \textbf{Mutable default state in classes.} \code{class C: items = []} shares the list across instances (it's class-level). Use \code{__init__}.
\end{itemize}
\end{minipage}
\end{construct}

\end{document}
